\documentclass[../main.tex]{subfiles}
\begin{document}
    \part{Expressions littérales}
        \subpart{Définition}
            \definition{expression_litterale}
            \begin{exemple}
                \begin{enumerate}[i),font=\bfseries]
                    \begin{multicols}{2}    
                        \item $2\times(6+b)$
                        \item $x\times 7$
                        \columnbreak
                        \item $2\times\pi\times r$
                        \item $(b\times h)\div 2$
                    \end{multicols}
                \end{enumerate}
            \end{exemple}
        \subpart{Conventions d'écriture}
            \regle{calcul_litteral_conventions_ecriture}
            \begin{exemple}
                \begin{enumerate}[i),font=\bfseries]
                    \begin{multicols}{2}
                        \item $4\times a$ se note $4a$.
                        \columnbreak
                        \item $3\times(y+8)$ se note $3(y+8)$.
                    \end{multicols}
                \end{enumerate}
            \end{exemple}
            \begin{remarque}\smallskip\par\noindent 
                \begin{enumerate}[a),font=\bfseries]
                    \begin{multicols}{2}
                        \item $1\times d$ se note $d$.
                        \item $(-1)\times x$ se note $-x$.
                        \columnbreak
                        \item $r\times r$ se note $r^2$.
                        \item $\pi\times\pi\times\pi$ se note $\pi^3$.
                    \end{multicols}
                \end{enumerate}
            \end{remarque}
    \part{Évaluer d'une expression littérale}
        \propriete{evaluer_une_expression_litterale}
        \begin{exemple}
            On évalue l'expression littérale $4a-12$ pour $a=3$, $a=0$ et $a=-2$.
            \vskip 10 pt
            Pour $a=3$ :
            $$\begin{array}{cccc}
                &4a&-&12\\
                =&4\times a&-&12\\
                =&4\times 3&-&12\\
                =&12&-&12\\
                =&&0&
            \end{array}$$
            \vskip 5 pt
            Pour $a=0$ :
            $$\begin{array}{cccc}
                &4a&-&12\\
                =&4\times a&-&12\\
                =&4\times 0&-&12\\
                =&0&-&12\\
                =&&-12&
            \end{array}$$
            \vskip 5 pt
            Pour $a=-2$ :
            $$\begin{array}{cccc}
                &4a&-&12\\
                =&4\times a&-&12\\
                =&4\times(-2)&-&12\\
                =&-8&-&12\\
                =&&-20&
            \end{array}$$
        \end{exemple}
    \part{Test d'égalité}
        \definition{egalite}
        \begin{exemple}
            \begin{enumerate}[i),font=\bfseries]
                \item L'égalité $0=1$ est {\bfseries fausse}.
                \item L'égalité $2\times 4=2^3$ est {\bfseries vraie}.
            \end{enumerate}
        \end{exemple}
        \methode{tester_une_egalite}
        \begin{exemple}
            \begin{enumerate}[i), font=\bfseries]
                \item Testons l'égalité $3x+2=5x+4$ pour $x=3$.
                \vskip 10pt
                \begin{tblr}{Q[wd=.45\linewidth]|Q[wd=.45\linewidth]}
                    On évalue le membre de gauche pour $x=3$ :
                    \par\vskip 10pt\hskip 30pt
                    $\begin{array}{cccc}
                        &3x&+&2\\
                        =&3\times x&+&2\\
                        =&4\times 3&+&2\\
                        =&12&+&2\\
                        =&&14&
                    \end{array}$
                    &
                    On évalue le membre de droite pour $x=3$ :
                    \par\vskip 10pt\hskip 30pt
                    $\begin{array}{cccc}
                        &5x&+&4\\
                        =&5\times x&+&4\\
                        =&5\times 3&+&4\\
                        =&15&+&4\\
                        =&&19&
                    \end{array}$
                \end{tblr}\vskip 10pt%
                Les deux membres de l'égalité sont différents donc l'égalité
                est fausse pour $x=3$.
                \item Testons l'égalité $6x=2x-8$ pour $x=-2$.
                \vskip 10pt
                \begin{tblr}{Q[wd=.45\linewidth]|Q[wd=.45\linewidth]}
                    On évalue le membre de gauche pour $x=-2$ :
                    \par\vskip 10pt\hskip 30pt
                    $\begin{array}{cccc}
                        &&6x&\\
                        =&6&\times&x\\
                        =&6&\times&(-2)\\
                        =&&-12&
                    \end{array}$
                    &
                    On évalue le membre de droite pour $x=-2$ :
                    \par\vskip 10pt\hskip 30pt
                    $\begin{array}{cccc}
                        &2x&-&8\\
                        =&2\times x&-&8\\
                        =&2\times(-2)&-&8\\
                        =&-4&-&8\\
                        =&&-12&
                    \end{array}$
                \end{tblr}\vskip 10pt%
                Les deux membres de l'égalité sont égaux donc l'égalité
                est vraie pour $x=-2$.
            \end{enumerate}
        \end{exemple}
        \newpage
        \part{Réduction d'une expression littérale}
        \definition{reduire_une_expression_litterale}
        \begin{exemple}
            \begin{enumerate}[i),font=\bfseries]
                \item$$\begin{array}{cccc}
                    &a&+&3+4\\
                    =&a&+&7
                \end{array}$$
                \item$$\begin{array}{cccccc}
                    &-2&+&b&+&2\\
                    =&b&+&-2&+&2\\
                    =&b&+&&0&\\
                    =&b&&&&
                \end{array}$$
                \item$$\begin{array}{cccccc}
                    &4&\times&3&\times&x\\
                    =&&12&&\times&x\\
                    =&&&&12x&
                \end{array}$$
            \end{enumerate}
        \end{exemple}
\end{document}